\documentclass[12pt,a4paper]{article}

\usepackage[a4paper,margin=2cm]{geometry}
\usepackage{fontspec}
\usepackage{polyglossia}
\setdefaultlanguage{russian}
\setmainfont{Times New Roman}
\newfontfamily\cyrillicfonttt{Courier New}
\usepackage{amsmath,amssymb}
\usepackage{booktabs}
\usepackage{longtable}
\usepackage{array}
\usepackage{enumitem}
\usepackage{indentfirst}
\usepackage[hidelinks]{hyperref}

\setlength{\parindent}{1.25cm}
\setlength{\parskip}{0.3em}
\setlist{nosep,leftmargin=1.5cm}
\emergencystretch=3em

\begin{document}

\begin{titlepage}
    \centering

    {\small Министерство науки и высшего образования Российской Федерации\par}
    {\small Федеральное государственное автономное образовательное учреждение высшего образования\par}
    {\small «Национальный исследовательский университет ИТМО»\par}

    \vspace{3cm}

    {\Large\bfseries Проект\par}
    \vspace{0.5cm}
    {\large по дисциплине «Физика»\par}
    \vspace{0.8cm}
    {\Large\bfseries Программный комплекс для численного моделирования распространения сигнала сотовой связи в городской среде\par}

    \vfill

    \begin{flushright}
        \begin{tabular}{rl}
            Выполнил: & Ларионов Владислав Васильевич \\
            Выполнил: & Кузьмин Дмитрий Анатольевич \\
            Преподаватель: & Сорокина Елена Константиновна
        \end{tabular}
    \end{flushright}

    \vfill

    {\large Санкт-Петербург, 2026\par}
\end{titlepage}

\tableofcontents
\newpage

\section{Назначение проекта}

Цель проекта --- построить интерактивную численную модель распространения радиосигнала в городской среде. Модель загружает реальный фрагмент города из OpenStreetMap, строит трехмерную сцену, размещает базовые станции на периметре расчетной области и показывает, как меняется качество связи при перемещении вышек, движении крупного объекта и появлении источника помехи.

Основной результат расчета --- трехмерная voxel-сетка качества связи. Для каждой ячейки определяется суммарная мощность полезного сигнала, мощность внешней помехи и отношение сигнал/шум. В интерфейсе отображается горизонтальный срез этой сетки на выбранной пользователем высоте.

Проект реализован как настольное Python-приложение на PySide6, PyVista и pyvistaqt. Пользователь может вращать трехмерную сцену, перемещать вышки, включать и отключать отдельные источники сигнала, запускать БПЛА, выбирать режим визуализации и пересчитывать карту покрытия.

\section{Структура реализации}

Код проекта разделен на несколько модулей:

\begin{itemize}
    \item \path{main.py} --- графический интерфейс, обработчики событий, фоновые задачи пересчета покрытия;
    \item \path{city/osm_parser.py} --- чтение OSM XML, извлечение зданий и метаданных, структура данных сцены;
    \item \path{city/projection.py} --- перевод географических координат в локальные метры;
    \item \path{city/scene.py} --- подготовка расчетной области, размещение вышек, добавление объектов в PyVista;
    \item \path{city/builder.py} --- построение геометрических мешей зданий, вышек, парома, БПЛА, волн и среза покрытия;
    \item \path{city/radio.py} --- трассировка радиолучей, расчет потерь, отражений, помехи и voxel-карты покрытия.
\end{itemize}

В качестве основного района используется файл \path{districts/willis_mega.osm}. Также есть уменьшенный район \path{districts/willis_mini.osm}, удобный для быстрых проверок и отладки.

\section{Геометрическая модель сцены}

\subsection{Переход к локальным координатам}

Исходные данные OpenStreetMap задаются широтой и долготой. Для небольшой городской области используется эквидистантная локальная проекция относительно опорной точки:

\[
x = R(\lambda - \lambda_0)\cos\varphi_0,
\qquad
y = R(\varphi - \varphi_0),
\]

где $R$ --- радиус Земли, $\varphi$ и $\lambda$ --- широта и долгота точки в радианах, $\varphi_0$ и $\lambda_0$ --- координаты опорной точки. Ось $x$ направлена на восток, ось $y$ --- на север. Для рассматриваемой области размером в сотни метров искажения такой проекции малы по сравнению с шагом расчетной сетки.

\subsection{Здания и расчетная область}

Из OSM извлекаются замкнутые контуры с тегом \texttt{building}. Если задана высота здания, она используется напрямую. Если высоты нет, применяется число этажей и средняя высота этажа:

\[
h = 3N_{\text{levels}}.
\]

Если не задана ни высота, ни число этажей, используется высота по умолчанию $9$ м.

Каждый контур здания вытягивается по вертикали в призму. Контуры площадью меньше $35$ м$^2$ отбрасываются как несущественные для визуализации и трассировки. Для размещения движущегося объекта из района удаляется узкая полоса зданий, имитирующая свободный водный или транспортный коридор.

Вокруг района строится прямоугольный контур и невидимые граничные стены высотой $400$ м. Эти стены участвуют в отражениях радиолучей и не дают лучам бесконечно уходить за пределы расчетной области.

\subsection{Базовые станции}

В текущей версии две базовые станции создаются автоматически на противоположных сторонах периметра. Их положение задается одним параметром $p \in [0, 1)$: первая вышка ставится в точку периметра с параметром $p$, вторая --- в точку $p + 0.5$. Это сохраняет симметричное расположение источников относительно области.

Источник радиосигнала находится не у основания вышки, а на высоте $90\%$ от ее полной высоты:

\[
T = (x_t, y_t, 0.9h_t),
\]

где $(x_t, y_t)$ --- центр основания вышки, $h_t$ --- высота вышки.

\section{Физическая модель радиосигнала}

\subsection{Расстояние и потери на распространение}

Для каждой точки расчета используется расстояние от источника сигнала до центра ячейки:

\[
d = \sqrt{(x - x_t)^2 + (y - y_t)^2 + (z - z_t)^2}.
\]

После отражения учитывается не только длина текущего участка, но и полный уже пройденный путь луча:

\[
d_{\text{cell}} = d_{\text{prev}} + d_{\text{seg}},
\]

где $d_{\text{prev}}$ --- путь до начала текущего участка, $d_{\text{seg}}$ --- расстояние от начала участка до центра ячейки.

Потери на распространение считаются по логарифмической модели:

\[
L(d) = L(d_0) + 10n\log_{10}\left(\frac{d}{d_0}\right).
\]

Такая модель выбрана потому, что она хорошо подходит для приближенной оценки затухания радиосигнала в городской среде без полного решения уравнений Максвелла. Параметр $n$ позволяет учесть более быстрое падение мощности в плотной застройке по сравнению со свободным пространством, а форма зависимости остается достаточно простой для интерактивного пересчета покрытия.

В реализации используются параметры:

\[
L(d_0) = 32\text{ дБ}, \qquad d_0 = 1\text{ м}, \qquad n = 2.6.
\]

Если расстояние меньше $d_0$, в формуле используется $d_0$, чтобы избежать некорректного поведения логарифма около нуля.

\subsection{Баланс принимаемой мощности}

Мощность принятого полезного сигнала от отдельного луча вычисляется как:

\[
P_r = P_t + G_t + G_r - L(d) - L_{\text{refl}},
\]

где $P_t$ --- мощность передатчика, $G_t$ --- усиление передающей антенны, $G_r$ --- усиление приемной антенны, $L(d)$ --- потери на распространение, $L_{\text{refl}}$ --- накопленные потери после отражений.

В программе заданы значения:

\[
P_t = 46\text{ дБм}, \qquad G_t = 15\text{ дБ}, \qquad G_r = 0\text{ дБ}.
\]

Начальный энергетический уровень до учета потерь равен $61$ дБм.

\subsection{Параметры модели}

Основные численные параметры, используемые в расчете, приведены в таблице.

\begin{center}
\begin{tabular}{lll}
\toprule
Параметр & Значение & Смысл \\
\midrule
$P_t$ & $46$ дБм & мощность базовой станции \\
$G_t$ & $15$ дБ & усиление передатчика \\
$G_r$ & $0$ дБ & усиление приемника \\
$L(d_0)$ & $32$ дБ & потери на опорном расстоянии \\
$d_0$ & $1$ м & опорное расстояние \\
$n$ & $2.6$ & показатель затухания \\
$P_{\text{noise}}$ & $-100$ дБм & уровень фонового шума \\
$P_{\text{uav}}$ & $20$ дБм & мощность помехи слабого БПЛА \\
$a$ & $15$ м & размер voxel-ячейки по умолчанию \\
$L_{\text{building}}$ & $8$ дБ & отражение от здания \\
$L_{\text{ferry}}$ & $3$ дБ & отражение от парома \\
$L_{\text{boundary}}$ & $14$ дБ & отражение от границы \\
\bottomrule
\end{tabular}
\end{center}

\subsection{Отражения}

Лучи отражаются от зданий, граничных стен и движущегося парома. Направление отраженного луча задается законом зеркального отражения:

\[
\vec r = \vec d - 2(\vec d \cdot \vec n)\vec n,
\]

где $\vec d$ --- направление падающего луча, $\vec n$ --- единичная нормаль поверхности, $\vec r$ --- направление после отражения.

При каждом отражении к дальнейшему пути добавляется фиксированная потеря:

\[
L_{\text{building}} = 8\text{ дБ}, \qquad
L_{\text{ferry}} = 3\text{ дБ}, \qquad
L_{\text{boundary}} = 14\text{ дБ}.
\]

Если тип поверхности не определен, используется потеря $10$ дБ. После нескольких отражений суммарная потеря равна:

\[
L_{\text{refl}} = \sum_{j=1}^{M} L_j.
\]

Количество отражений ограничено параметром интерфейса. По умолчанию используется два отражения.

\subsection{Дискретная трассировка лучей}

От каждой активной базовой станции выпускается конечное число лучей. Направления распределяются по сфере приближенно равномерно с помощью решетки Фибоначчи. Такой подход дает не аналитическое поле, а численное приближение: ячейка получает вклад только тогда, когда через нее прошел один или несколько расчетных лучей.

Для каждого прямого участка луча программа выбирает ячейки voxel-сетки, пересеченные этим участком, и добавляет в них вклад мощности. Слишком слабые участки ниже $-130$ дБм отбрасываются.

Вклады нескольких лучей суммируются в линейной шкале мощности:

\[
P_{\text{mW}} = 10^{P_{\text{dBm}}/10},
\]

\[
P_{\text{signal}} = \sum_{m=1}^{M}P_m.
\]

Сложение выполняется именно в мВт, так как значения в дБм нельзя складывать напрямую.

\section{Карта покрытия и SNR}

Расчетная область разбивается на кубические ячейки с шагом $a$. По умолчанию:

\[
a = 15\text{ м}.
\]

Размеры сетки:

\[
N_x = \left\lceil \frac{x_{\max} - x_{\min}}{a} \right\rceil,
\qquad
N_y = \left\lceil \frac{y_{\max} - y_{\min}}{a} \right\rceil,
\qquad
N_z = \left\lceil \frac{z_{\max} - z_{\min}}{a} \right\rceil.
\]

Качество связи отображается через отношение сигнал/шум:

\[
\mathrm{SNR} =
\frac{P_{\text{signal}}}
{P_{\text{noise}} + P_{\text{jam}}},
\]

где $P_{\text{signal}}$ --- суммарная мощность полезного сигнала, $P_{\text{noise}}$ --- фоновый шум, $P_{\text{jam}}$ --- внешняя помеха от БПЛА. Фоновый шум принят равным:

\[
P_{\text{noise}} = -100\text{ дБм}.
\]

Для визуализации используется значение в децибелах:

\[
\mathrm{SNR}_{\text{dB}} =
10\log_{10}
\left(
\frac{P_{\text{signal}}}
{P_{\text{noise}} + P_{\text{jam}}}
\right).
\]

Пользователь выбирает высоту горизонтального среза. Если высота находится вне расчетной сетки, она ограничивается допустимым диапазоном.

\section{Динамические объекты}

\subsection{Паром}

В сцене присутствует движущийся крупный объект, моделирующий паром или лайнер. Он движется вдоль выделенного коридора, отображается как составной трехмерный меш и участвует в отражениях радиолучей. При включенном пароме его положение входит в ключ кэша покрытия, поэтому разные положения объекта дают разные результаты расчета.

Для ускорения используются статические траектории лучей без учета парома. Если статический луч заведомо не пересекает область парома, он переиспользуется. Если пересечение возможно, луч пересчитывается уже с полным набором отражателей.

\subsection{Режимы БПЛА}

В интерфейсе реализованы два режима БПЛА. Оба режима используют одинаковую геометрию, высоту и траекторию: аппарат летит поперек расчетной области за $10$ секунд. Его координаты зависят от текущего прогресса полета:

\[
U = (x_u, y_u, z_u).
\]

Слабый БПЛА является источником внешней помехи. Для каждой ячейки вычисляется расстояние до БПЛА:

\[
d_{\text{uav}} =
\sqrt{(x - x_u)^2 + (y - y_u)^2 + (z - z_u)^2}.
\]

Мощность помехи рассчитывается по той же логарифмической модели потерь:

\[
P_{\text{jam}} = P_{\text{uav}} - L(d_{\text{uav}}).
\]

В реализации:

\[
P_{\text{uav}} = 20\text{ дБм}, \qquad z_u = 120\text{ м}.
\]

После расчета мощность помехи переводится в мВт и добавляется к знаменателю SNR. Если слабый БПЛА не активен, внешняя помеха от него равна нулю.

Сильный БПЛА не добавляет шум в знаменатель SNR. Вместо этого на каждом шаге полета вычисляется горизонтальное расстояние от БПЛА до активных вышек. Все активные вышки в радиусе $320$ м временно отключаются: их состояние переводится в неактивное, в сцене они становятся красными, а соответствующий тумблер в интерфейсе снимается. Когда БПЛА выходит из зоны воздействия или завершает полет, автоматически включаются только те вышки, которые были отключены именно сильным БПЛА. Пользовательские отключения при этом не восстанавливаются.

\section{Режимы визуализации}

В интерфейсе доступны три режима:

\begin{enumerate}
    \item динамические волны от активных вышек;
    \item трассированные радиолучи с цветом по мощности;
    \item горизонтальный срез карты покрытия с цветом по SNR.
\end{enumerate}

Динамические волны являются визуальной демонстрацией распространения фронта сигнала от вышек. Они не используются напрямую в численном расчете покрытия: расчет выполняется трассировкой лучей и суммированием мощностей в voxel-сетке.

\section{Кэширование и производительность}

Расчет покрытия является самой тяжелой частью проекта, поэтому результаты сохраняются в памяти и на диске. Ключ кэша включает:

\begin{itemize}
    \item размер voxel-ячейки;
    \item число лучей;
    \item число отражений;
    \item включение и положение парома;
    \item состояние и положение слабого БПЛА;
    \item положение, высоту и состояние вышек, включая отключение сильным БПЛА;
    \item геометрию зданий и границу расчетной области.
\end{itemize}

Файловый кэш хранится в директории \path{.cache/coverage}. Он не является исходными данными проекта и может быть удален: при необходимости приложение пересчитает покрытие заново.

\section{Что реализовано}

В текущей версии реализованы:

\begin{enumerate}
    \item загрузка OSM-данных района Чикаго;
    \item построение трехмерной городской сцены;
    \item автоматическое размещение двух базовых станций на периметре;
    \item интерактивное перемещение вышек с сохранением симметрии;
    \item включение и отключение каждой вышки;
    \item движущийся паром, участвующий в отражениях;
    \item запуск слабого БПЛА, создающего внешнюю помеху;
    \item запуск сильного БПЛА, временно отключающего ближайшие вышки;
    \item визуализация динамических волн;
    \item трассировка радиолучей с отражениями;
    \item расчет мощности полезного сигнала в voxel-сетке;
    \item расчет SNR с учетом фонового шума и помехи от БПЛА;
    \item отображение горизонтального среза качества покрытия;
    \item кэширование расчетов в памяти и на диске.
\end{enumerate}

\section{Ограничения модели}

Модель является инженерным приближением и имеет несколько ограничений:

\begin{itemize}
    \item не рассчитываются фазы электромагнитных волн и интерференционная картина;
    \item вклад лучей суммируется как мощность, а не как комплексные амплитуды поля;
    \item материалы зданий не задаются индивидуально, потери при отражениях фиксированы по типу объекта;
    \item дифракция вокруг углов зданий не моделируется;
    \item высоты части зданий берутся из OSM, а при отсутствии данных используются приближенные значения;
    \item точность карты зависит от числа трассируемых лучей и размера voxel-ячейки.
\end{itemize}

Эти упрощения позволяют получить интерактивную скорость работы и наглядную карту качества связи, сохраняя связь с базовыми физическими законами: геометрическим расстоянием, логарифмическими потерями, отражением лучей и отношением сигнал/шум.

\section{Вывод}

В результате создан рабочий программный прототип для исследования распространения радиосигнала в городской среде. Проект объединяет реальные геоданные, трехмерную геометрию, численную трассировку лучей и расчет SNR. Пользователь может наблюдать, как меняется покрытие при изменении положения базовых станций, движении крупного объекта и появлении БПЛА как источника помехи.

Текущая реализация завершает поставленную задачу на уровне интерактивной инженерной модели: она строит городскую сцену, рассчитывает распространение сигнала по заданной физической модели и позволяет сравнивать качество покрытия при разных состояниях системы. Принятые упрощения явно отделены от реализованной части модели и не мешают использовать проект для демонстрации влияния городской геометрии, отражений, движущихся объектов и внешней помехи на качество связи.

\end{document}
